\documentclass[11pt,a4paper]{article}

% Template visual Matriz Aprovação. Compile com XeLaTeX ou LuaLaTeX.
\usepackage{fontspec}
\usepackage{polyglossia}
\setmainlanguage{portuguese}
\IfFontExistsTF{Aptos}{
  \setmainfont{Aptos}
  \setsansfont{Aptos}
}{
  \IfFontExistsTF{Calibri}{
    \setmainfont{Calibri}
    \setsansfont{Calibri}
  }{
    \IfFontExistsTF{TeX Gyre Heros}{
      \setmainfont{TeX Gyre Heros}
      \setsansfont{TeX Gyre Heros}
    }{
      \setmainfont{Latin Modern Sans}
      \setsansfont{Latin Modern Sans}
    }
  }
}
\usepackage[a4paper,margin=2cm,headheight=16pt]{geometry}
\usepackage{graphicx}
\usepackage{xcolor}
\usepackage{tikz}
\usepackage{tcolorbox}
\usepackage{enumitem}
\usepackage{booktabs}
\usepackage{tabularx}
\usepackage{amssymb}
\usepackage{fancyhdr}
\usepackage{titlesec}
\usepackage{hyperref}

\definecolor{matrizlime}{HTML}{CBFF4D}
\definecolor{matrizink}{HTML}{101313}
\definecolor{matrizmuted}{HTML}{596363}
\definecolor{matrizline}{HTML}{DDE3DF}
\definecolor{matrizsoft}{HTML}{F3F7EC}

\hypersetup{colorlinks=true,linkcolor=matrizink,urlcolor=matrizink}
\pagestyle{fancy}
\fancyhf{}
\lhead{\textsf{\small MATRIZ APROVAÇÃO}}
\rhead{\textsf{\small APOSTILA}}
\cfoot{\textsf{\small \thepage}}
\renewcommand{\headrulewidth}{0.4pt}
\renewcommand{\headrule}{\hbox to\headwidth{\color{matrizline}\leaders\hrule height \headrulewidth\hfill}}

\titleformat{\section}{\Large\bfseries\color{matrizink}}{\thesection}{0.6em}{}
\titleformat{\subsection}{\large\bfseries\color{matrizink}}{\thesubsection}{0.6em}{}
\setlist{nosep,leftmargin=*}

\newtcolorbox{foco}{colback=matrizsoft,colframe=matrizlime,boxrule=1pt,arc=3pt,left=10pt,right=10pt,top=8pt,bottom=8pt}
\newtcolorbox{lei}{colback=matrizink,colframe=matrizink,coltext=white,boxrule=0pt,arc=3pt,left=10pt,right=10pt,top=8pt,bottom=8pt}
\newtcolorbox{alerta}{colback=white,colframe=matrizline,boxrule=0.6pt,arc=3pt,left=10pt,right=10pt,top=8pt,bottom=8pt}

% Logo usada na capa. Caminho relativo à pasta onde o .tex é compilado
% (materiais/<disciplina>/). Ajuste se a estrutura mudar.
\newcommand{\logomatriz}{../../public/logo.png}

\begin{document}

\begin{titlepage}
  \thispagestyle{empty}
  \begin{center}
    \includegraphics[width=0.55\linewidth]{\logomatriz}\par
  \end{center}
  \vspace{2.2cm}
  {\sffamily\fontsize{28}{32}\selectfont\bfseries Apostila — Redação: Texto dissertativo-argumentativo para o ENEM\par}
  \vspace{0.4cm}
  {\sffamily\large ENEM\par}
  \vfill
  \begin{foco}
    \textsf{\textbf{Foco da apostila}}\\[3pt]
    cada parágrafo precisa cumprir uma função. A introdução delimita o problema e apresenta a tese; os desenvolvimentos provam a tese; a conclusão retoma o raciocínio e detalha a intervenção.
  \end{foco}
  \vspace{0.7cm}
  {\sffamily\small Atualizado em 16 de agosto de 2026\quad ·\quad INEP\par}
  \vspace{1cm}
  {\sffamily\small matrizaprova.com\par}
\end{titlepage}

\tableofcontents
\newpage

% Conteúdo gerado entra aqui.
\section*{O que cai}

A redação do ENEM exige um texto dissertativo-argumentativo em modalidade escrita formal da língua portuguesa. O participante deve defender um ponto de vista sobre um tema de ordem social, científica, cultural ou política e apresentar uma proposta de intervenção para o problema discutido.

Os pontos decisivos são:

\begin{enumerate}
  \item compreender exatamente o recorte temático e não escrever sobre um assunto
\end{enumerate}

   apenas relacionado;

\begin{enumerate}
  \item apresentar tese clara e argumentos desenvolvidos;
  \item organizar introdução, desenvolvimento e conclusão com progressão lógica;
  \item usar repertório sociocultural pertinente e articulado ao argumento;
  \item propor intervenção detalhada, viável e compatível com os direitos humanos.
\end{enumerate}

\textbf{Regra de prova:} cada parágrafo precisa cumprir uma função. A introdução delimita o problema e apresenta a tese; os desenvolvimentos provam a tese; a conclusão retoma o raciocínio e detalha a intervenção.

\section*{Teoria essencial}

\subsection*{1. Entender o tema e o recorte}

Antes de escrever, transforme o tema em uma pergunta-problema. O tema pode apresentar uma situação, um grupo afetado, uma causa ou uma consequência. Todos esses elementos formam o recorte que precisa aparecer no texto.

Faça três marcações:

\begin{enumerate}
  \item \textbf{assunto geral:} sobre o que se fala;
  \item \textbf{problema específico:} qual dificuldade deve ser discutida;
  \item \textbf{comando implícito:} o texto deve explicar causas, consequências,
\end{enumerate}

   desafios ou formas de enfrentamento.

Se o tema for “desafios para combater um problema no Brasil”, não basta definir o problema nem escrever uma lista de soluções. É preciso explicar por que ele persiste e como pode ser enfrentado no contexto brasileiro.

O tangenciamento ocorre quando o texto aborda apenas uma parte secundária do tema. Para evitar isso, repita mentalmente o recorte em cada etapa: a tese, os argumentos e a proposta precisam responder ao mesmo problema.

\subsection*{2. Projeto de texto}

O projeto de texto é o planejamento dos argumentos antes da redação. Em cinco minutos, escreva uma tese e duas causas ou dimensões do problema.

Um modelo funcional:

\begin{itemize}
  \item tese: o problema persiste por causa de A e B;
  \item desenvolvimento 1: explique A, suas origens e seus efeitos;
  \item desenvolvimento 2: explique B, suas origens e seus efeitos;
  \item conclusão: enfrente A e B com uma proposta articulada.
\end{itemize}

Os argumentos podem ser institucionais, sociais, econômicos, culturais, históricos ou tecnológicos. O importante não é acumular causas, mas desenvolver duas ideias com relação clara com a tese.

Evite escrever a introdução inteira antes de saber o que será defendido. Uma tese genérica, como “é preciso resolver esse problema”, não orienta os parágrafos. Uma tese eficiente mostra a direção do texto: “A persistência do problema decorre da insuficiência de políticas públicas e da naturalização de determinadas práticas sociais”.

\subsection*{3. Introdução}

A introdução pode ter três movimentos:

\begin{enumerate}
  \item contextualização breve;
  \item apresentação do problema;
  \item tese com os dois eixos argumentativos.
\end{enumerate}

A contextualização não precisa ser uma citação famosa. Pode ser uma observação histórica, uma definição, um contraste ou uma referência a princípio constitucional, desde que não substitua a discussão do tema.

Modelo:

\begin{alerta}
Embora [contexto ou avanço], [problema] ainda persiste no Brasil. Nesse cenário, a dificuldade é sustentada por [argumento 1] e [argumento 2], fatores que exigem enfrentamento articulado.
\end{alerta}

O modelo não deve ser copiado mecanicamente. Ele serve para conferir se a introdução contém contexto, problema e tese. A tese pode aparecer em uma ou duas frases, desde que os argumentos estejam identificáveis.

\subsection*{4. Desenvolvimento argumentativo}

Cada desenvolvimento deve apresentar uma ideia central e provar essa ideia. Uma estrutura segura é:

\begin{enumerate}
  \item tópico frasal: afirmação do argumento;
  \item explicação: como o fator produz ou agrava o problema;
  \item repertório ou exemplo: referência que sustenta a análise;
  \item consequência: efeito concreto sobre a sociedade;
  \item fechamento: ligação com a tese.
\end{enumerate}

Não basta citar uma instituição, um filósofo ou uma lei. O repertório precisa ser produtivo, isto é, explicar o argumento. A frase “Segundo determinado autor, a sociedade é desigual” só ajuda se o texto mostrar como essa desigualdade se relaciona ao problema proposto.

Um parágrafo de desenvolvimento não é uma coleção de frases independentes. Use conectivos para indicar causa, consequência, oposição e conclusão: “desse modo”, “além disso”, “por conseguinte”, “contudo”, “portanto”. O conectivo deve corresponder à relação lógica real.

\subsection*{5. Repertório sociocultural}

Repertório é conhecimento externo pertinente ao tema. Pode vir de História, Sociologia, Filosofia, Literatura, legislação, pesquisas, obras artísticas, conceitos científicos ou acontecimentos conhecidos.

Um repertório forte possui três características:

\begin{itemize}
  \item pertinência: relaciona-se diretamente com o argumento;
  \item legitimidade: é reconhecível e não inventado;
  \item produtividade: é explicado e usado para defender a tese.
\end{itemize}

Não invente dados, autores, títulos ou frases. Se não tiver certeza da autoria de uma citação, use a ideia sem atribuição literal ou escolha outro repertório. Informação falsa pode comprometer a credibilidade e a argumentação.

É melhor usar dois repertórios bem explicados do que cinco menções decorativas. Uma referência constitucional, por exemplo, deve ser ligada ao dever do Estado, à proteção de um direito ou à cidadania discutida no tema.

\subsection*{6. Coesão e coerência}

Coerência é a lógica global do texto. Coesão é a ligação linguística entre suas partes. O texto deve avançar sem contradizer a tese ou repetir a mesma ideia.

Use recursos de coesão:

\begin{itemize}
  \item pronomes e expressões de retomada: “esse cenário”, “tal medida”;
  \item substituição lexical: “problema”, “desafio”, “situação”;
  \item conectivos: “porque”, “assim”, “entretanto”, “além disso”;
  \item paralelismo: estruturas semelhantes para ideias semelhantes.
\end{itemize}

Evite iniciar todos os parágrafos com “é importante”. Também evite “porém” como conectivo de qualquer relação. Se o segundo período explicar o primeiro, use “pois” ou “porque”; se apresentar oposição, use “porém” ou “contudo”.

\subsection*{7. Conclusão e proposta de intervenção}

A conclusão deve retomar o problema sem repetir a introdução e apresentar uma proposta de intervenção. Uma proposta completa contém cinco elementos:

\begin{enumerate}
  \item \textbf{agente:} quem executará a medida;
  \item \textbf{ação:} o que será feito;
  \item \textbf{meio ou modo:} como a ação será realizada;
  \item \textbf{finalidade:} qual resultado se pretende alcançar;
  \item \textbf{detalhamento:} especificação de um agente, ação ou efeito.
\end{enumerate}

Exemplo estrutural:

\begin{alerta}
Portanto, o Ministério responsável deve implementar [ação], por meio de [modo], a fim de [finalidade]. Paralelamente, [outro agente] deve [segunda ação], com [detalhamento], para reduzir [efeito do problema].
\end{alerta}

O agente deve ter competência para executar a medida. “A sociedade deve resolver” é vago; “o Ministério da Educação, em parceria com as redes de ensino” é mais preciso quando o tema envolver educação. A proposta precisa respeitar os direitos humanos e não pode sugerir violência, censura ou retirada de direitos.

Quando o tema tiver duas causas, a intervenção deve enfrentá-las. Uma campanha de conscientização isolada pode ser insuficiente para um problema que também depende de fiscalização, orçamento ou acesso a serviços.

\subsection*{8. Linguagem formal e revisão}

A modalidade formal exige atenção a concordância, regência, pontuação, acentuação e escolha vocabular. Não é necessário usar palavras difíceis. Clareza e precisão valem mais que frases ornamentadas.

Na revisão, faça quatro leituras rápidas:

\begin{enumerate}
  \item tema: respondi exatamente ao recorte?
  \item estrutura: há tese, dois argumentos e intervenção?
  \item conexão: cada parágrafo se liga ao anterior?
  \item língua: corrigi concordância, pontuação, repetição e grafia?
\end{enumerate}

Não introduza uma ideia nova na última linha. A conclusão deve fechar o raciocínio, não abrir um terceiro desenvolvimento.

\section*{Como a banca cobra}

\begin{itemize}
  \item A redação pode perder desempenho quando discute o assunto amplo, mas não o recorte específico do tema.
  \item Citação sem explicação não substitui argumento.
  \item Proposta com apenas “conscientizar” costuma ser genérica se não houver agente, meio e finalidade.
  \item Conclusão que repete a tese sem intervenção fica incompleta.
  \item Frases longas com muitas orações aumentam o risco de falha de pontuação e concordância.
  \item Repertório decorado precisa ser adaptado ao tema; não force uma referência que não sustenta a ideia.
\end{itemize}

\subsection*{Exemplos resolvidos}

\subsubsection*{Exemplo 1}

Tema hipotético: desafios para ampliar a inclusão digital de idosos no Brasil.

Tese fraca: “A inclusão digital é importante e deve ser ampliada.”

Tese eficiente: “A inclusão digital de idosos é dificultada pela oferta insuficiente de formação tecnológica e pela falta de acessibilidade nos serviços digitais.”

A segunda tese apresenta dois eixos que podem organizar os desenvolvimentos e a intervenção. O primeiro parágrafo pode discutir formação; o segundo, design e acessibilidade.

\subsubsection*{Exemplo 2}

Proposta incompleta: “É preciso conscientizar a população sobre o problema.”

Proposta desenvolvida: “O Ministério responsável deve promover oficinas gratuitas de letramento digital em centros públicos, com instrutores e materiais adaptados, a fim de ampliar a autonomia dos idosos no uso de serviços on-line. Além disso, as plataformas públicas devem realizar testes de acessibilidade e oferecer atendimento presencial alternativo, para impedir que a digitalização exclua usuários com dificuldades.”

A versão desenvolvida apresenta agentes, ações, meios, finalidades e detalhamento. Também enfrenta os dois eixos da tese.

\section*{Quadro-resumo}

\begin{tabularx}{\linewidth}{lX}
\toprule
 \textbf{Parte} & \textbf{Pergunta que deve responder} \\
\midrule
 introdução & qual problema será discutido e qual é a tese? \\
 desenvolvimento 1 & qual é a primeira causa ou dimensão? \\
 desenvolvimento 2 & qual é a segunda causa ou dimensão? \\
 conclusão & como o problema será enfrentado? \\
 agente & quem executará a medida? \\
 ação & o que será feito? \\
 meio & como será feito? \\
 finalidade & para qual resultado? \\
 detalhamento & qual especificação torna a proposta concreta? \\
\bottomrule
\end{tabularx}


\section*{Questões de fixação}

\subsection*{Questão 1 — (questão inédita)}

Em um projeto de texto para o ENEM, a tese tem como função principal:

A) apresentar todos os dados que serão usados. B) resumir exclusivamente a proposta de intervenção. C) indicar a posição defendida e orientar os argumentos. D) substituir os parágrafos de desenvolvimento. E) repetir literalmente o tema da proposta.

\begin{alerta}
\textbf{Gabarito: C.} A tese apresenta o ponto de vista e os eixos argumentativos. Os dados e a intervenção serão desenvolvidos nas partes adequadas do texto.
\end{alerta}

\subsection*{Questão 2 — (questão inédita)}

Qual repertório é mais produtivo em uma redação sobre desigualdade de acesso à educação?

A) uma citação famosa sem explicação. B) um conceito educacional explicado e ligado ao problema. C) uma lista de autores não relacionados ao tema. D) um dado estatístico inventado para reforçar a tese. E) uma referência histórica usada apenas na conclusão.

\begin{alerta}
\textbf{Gabarito: B.} Repertório produtivo é pertinente, legítimo e articulado ao argumento. Citação decorativa ou informação inventada não cumpre essa função.
\end{alerta}

\subsection*{Questão 3 — (questão inédita)}

Uma proposta de intervenção completa deve apresentar, entre outros elementos:

A) apenas uma opinião pessoal. B) agente, ação, meio, finalidade e detalhamento. C) uma nova tese sem relação com o desenvolvimento. D) uma promessa de resultado imediato. E) somente uma campanha publicitária.

\begin{alerta}
\textbf{Gabarito: B.} Esses elementos tornam a proposta concreta, executável e relacionada ao problema discutido.
\end{alerta}

\subsection*{Questão 4 — (questão inédita)}

Em “A medida é necessária, portanto deve ser acompanhada de fiscalização”, o termo “portanto” estabelece relação de:

A) oposição. B) condição. C) conclusão. D) concessão. E) comparação.

\begin{alerta}
\textbf{Gabarito: C.} A segunda oração apresenta uma conclusão decorrente da afirmação anterior.
\end{alerta}

\subsection*{Questão 5 — (questão inédita)}

Para evitar o tangenciamento, o candidato deve:

A) abordar qualquer assunto relacionado ao tema. B) repetir o título em todos os parágrafos. C) responder ao recorte específico ao longo de todo o texto. D) evitar apresentar uma tese definida. E) concentrar-se apenas na proposta de intervenção.

\begin{alerta}
\textbf{Gabarito: C.} O recorte deve orientar a tese, os argumentos e a intervenção. Tratar apenas do assunto amplo pode afastar o texto da proposta.
\end{alerta}

\section*{Revisão final}

\begin{itemize}
  \item[$\square$] Identifiquei o assunto, o problema e o recorte do tema.
  \item[$\square$] Escrevi uma tese com dois eixos argumentativos.
  \item[$\square$] Desenvolvi uma causa ou dimensão em cada parágrafo.
  \item[$\square$] Expliquei o repertório em vez de apenas citá-lo.
  \item[$\square$] Usei conectivos de acordo com a relação lógica.
  \item[$\square$] Retomei o problema na conclusão.
  \item[$\square$] Indiquei agente, ação, meio, finalidade e detalhamento.
  \item[$\square$] Mantive a proposta compatível com os direitos humanos.
  \item[$\square$] Revisei língua formal, pontuação e repetição.
\end{itemize}

\end{document}
